% ============================================================
%  preamble.tex — Packages & Formatting Config
%  แม่แบบเอกสารโครงงาน สาขาวิชาวิทยาการคอมพิวเตอร์
% ============================================================

% ─── Encoding & Font ─────────────────────────────────────────
\usepackage{fontspec}
\usepackage{polyglossia}
\setmainlanguage{thai}
\setotherlanguage{english}
\setmainfont{TH Sarabun New}[
  BoldFont       = {TH Sarabun New Bold},
  ItalicFont     = {TH Sarabun New Italic},
  BoldItalicFont = {TH Sarabun New Bold Italic},
  Scale          = 1.0
]
\IfFontExistsTF{Courier New}{\setmonofont{Courier New}}{\setmonofont{Latin Modern Mono}}
\newfontfamily\thaifonttt{TH Sarabun New}

% ─── Thai Alphabetical Page Numbering ────────────────────────
% ใช้แทน roman สำหรับส่วนนำ: ก ข ค ง จ ฉ ช ซ ...
\newcommand{\thaialph}[1]{%
  \ifcase#1%
  \or ก\or ข\or ค\or ง\or จ%
  \or ฉ\or ช\or ซ\or ฌ\or ญ%
  \or ฎ\or ฏ\or ฐ\or ฑ\or ฒ%
  \or ณ\or ด\or ต\or ถ\or ท%
  \or ธ\or น\or บ\or ป\or ผ%
  \or ฝ\or พ\or ฟ\or ภ\or ม%
  \else ?\fi%
}

% ─── Page Layout ─────────────────────────────────────────────
\usepackage[
  a4paper,
  top        = 3.81cm,  % ขอบบน 1.5 นิ้ว (body เริ่มที่นี่ เพราะใช้ includehead)
  bottom     = 2.54cm,  % ขอบล่าง 1.0 นิ้ว
  left       = 3.81cm,  % ขอบซ้าย 1.5 นิ้ว
  right      = 2.54cm,  % ขอบขวา 1.0 นิ้ว
  headheight = 18pt,    % ความสูง header box (~0.635 ซม.)
  headsep    = 1.27cm,  % ระยะ header ถึง body
                        % → header bottom = top − headsep = 3.81 − 1.27 = 2.54 ซม. (1 นิ้ว) จากขอบบน ✓
                        % → body เริ่มที่ top = 3.81 ซม. (1.5 นิ้ว) จากขอบบน ✓
]{geometry}
% \usepackage{showframe}  % DEBUG: แสดงเส้นขอบ margin/header/body/footer

% ─── Line Spacing & Paragraph ────────────────────────────────
\usepackage{setspace}
\usepackage{indentfirst}   % ย่อหน้าแรกหลัง heading ก็ indent ด้วย
\onehalfspacing
\setlength{\parindent}{1.27cm}
\setlength{\parskip}{0pt}
% ตัดบรรทัดภาษาไทย (XeTeX ICU) → ชิดขอบขวาสม่ำเสมอ
\XeTeXlinebreaklocale "th"
\XeTeXlinebreakskip = 0pt plus 0.5pt minus 0.5pt
\XeTeXlinebreakpenalty = 0
% ยืดหยุ่นกับบรรทัดที่มีคำอังกฤษยาว
\tolerance=2000
\emergencystretch=1.5em
\hbadness=2000

% ─── Headers & Page Numbers ──────────────────────────────────
\usepackage{fancyhdr}
\pagestyle{fancy}
\fancyhf{}
\fancyhead[R]{\fontsize{14}{16}\selectfont\thepage}  % เลขหน้ามุมบนขวา 14pt ตาม KKU format
\renewcommand{\headrulewidth}{0pt}
% ไม่แสดงเลขหน้าในหน้าแรกของบท (plain style)
\fancypagestyle{plain}{
  \fancyhf{}
  \renewcommand{\headrulewidth}{0pt}
}

% ─── Continuation title for multi-page lists (TOC/LOT/LOF) ───
% ใส่ "(ต่อ)" ลงใน body ของหน้าถัดไป ไม่ใช่ header
\newcommand{\toccontlabel}{สารบัญ}
\newif\iflstcont \lstcontfalse

\newcommand{\contlistpage}{
  \afterpage{%
    \iflstcont
      \thispagestyle{fancy}%
      \begin{center}%
        \fontsize{18}{22}\selectfont\bfseries \toccontlabel\ (ต่อ)%
      \end{center}%
      \noindent\hspace*{0pt}\hfill{\normalfont\fontsize{14}{18}\selectfont หน้า}%
      \vspace{6pt}%
      \contlistpage
    \fi
  }%
}

% pagestyle สำหรับหน้าแนวนอน (Landscape)
% physical top-right → มุมล่างขวาเมื่ออ่านในแนวนอน ← ตาม KKU format: เลขหน้ามุมบนขวา
\fancypagestyle{lscaped}{
  \fancyhf{}
  \fancyhead[R]{\thepage}
  \renewcommand{\headrulewidth}{0pt}
  \renewcommand{\footrulewidth}{0pt}
}

% ─── Chapter / Section Formatting ────────────────────────────
\usepackage{titlesec}

% กว้าง box เลข section = \parindent เสมอ → ชื่อตรงกับย่อหน้า
\newlength{\sechangindent}
\setlength{\sechangindent}{\parindent}   % 1.27cm

% บทที่: 18pt หนา, จัดกึ่งกลาง — เลขบทและชื่อบทคนละบรรทัด
\titleformat{\chapter}[block]
  {\centering\fontsize{18}{22}\selectfont\bfseries}
  {บทที่\ \thechapter}{0pt}{\\\centering}
  [\global\setlength{\parindent}{\sechangindent}]
\titlespacing{\chapter}{0pt}{-\topskip}{18pt}

% Section: 16pt หนา
\titleformat{\section}[hang]
  {\fontsize{16}{20}\selectfont\bfseries}
  {\makebox[\sechangindent][l]{\thesection}}{0pt}{}
  [\global\setlength{\parindent}{\sechangindent}]
\titlespacing{\section}{0pt}{12pt}{6pt}

% Subsection: 14pt หนา
\titleformat{\subsection}[hang]
  {\fontsize{14}{18}\selectfont\bfseries}
  {\makebox[\sechangindent][l]{\thesubsection}}{0pt}{}
  [\global\setlength{\parindent}{\dimexpr2\sechangindent\relax}]
\titlespacing{\subsection}{\sechangindent}{10pt}{4pt}

% Subsection ไม่หนา — ใช้ \subsectionmd{ชื่อ} สำหรับกรณีพิเศษ
\newcommand{\subsectionmd}[1]{%
  \titleformat{\subsection}[hang]%
    {\fontsize{14}{18}\selectfont}%
    {\makebox[\sechangindent][l]{\thesubsection}}{0pt}{}%
    [\global\setlength{\parindent}{\dimexpr2\sechangindent\relax}]%
  \titlespacing{\subsection}{\sechangindent}{10pt}{4pt}%
  \subsection{#1}%
  \titleformat{\subsection}[hang]%
    {\fontsize{14}{18}\selectfont\bfseries}%
    {\makebox[\sechangindent][l]{\thesubsection}}{0pt}{}%
    [\global\setlength{\parindent}{\dimexpr2\sechangindent\relax}]%
  \titlespacing{\subsection}{\sechangindent}{10pt}{4pt}%
}

% Subsubsection: 14pt ปกติ — ใช้ format (1), (2), (3)
\setcounter{secnumdepth}{3}  % ให้ \subsubsection มีเลขลำดับด้วย
\renewcommand{\thesubsubsection}{(\arabic{subsubsection})}

% คำนวณระยะออฟเซตให้ (2.1) ตรงกับอักษรตัวแรกของข้อความระดับ 3
\newlength{\subsuboffset}
\settowidth{\subsuboffset}{\fontsize{14}{18}\selectfont (0)\hspace{1em}}

\titleformat{\subsubsection}[hang]
  {\fontsize{14}{18}\selectfont}
  {\thesubsubsection}{1em}{}
  [\global\setlength{\parindent}{\dimexpr2\sechangindent+\subsuboffset\relax}]
\titlespacing{\subsubsection}{\dimexpr2\sechangindent\relax}{8pt}{4pt}
% reset เลข subsubsection ทุกครั้งที่ขึ้น subsection ใหม่
\makeatletter
\@addtoreset{subsubsection}{subsection}
\makeatother

% Subsubsubsection (ระดับ 4): 14pt ปกติพิมพ์ (N.N)
% สำหรับผู้ใช้ที่ต้องการใช้ \subsubsubsection
\newcommand{\subsubsubsection}[1]{\paragraph{#1}}
\setcounter{secnumdepth}{4}  % ให้ \paragraph มีเลขลำดับระดับ 4
\renewcommand{\theparagraph}{(\arabic{subsubsection}.\arabic{paragraph})}
\titleformat{\paragraph}[runin]
  {\fontsize{14}{18}\selectfont}
  {\theparagraph}{1em}{}
\titlespacing{\paragraph}{\dimexpr2\sechangindent+\subsuboffset\relax}{8pt}{4pt}
\makeatletter
\@addtoreset{paragraph}{subsubsection}
\makeatother

% Subsubsection ไม่มีลำดับ — ใช้ \subsubsectionmd{ชื่อ} สำหรับกรณีพิเศษ
\newcommand{\subsubsectionmd}[1]{%
  \titleformat{\subsubsection}[hang]%
    {\fontsize{14}{18}\selectfont}%
    {}{0pt}{}%
    [\global\setlength{\parindent}{\dimexpr2.675\sechangindent\relax}]%
  \titlespacing{\subsubsection}{\dimexpr2\sechangindent\relax}{8pt}{4pt}%
  \subsubsection*{#1}%
  \titleformat{\subsubsection}[hang]%
    {\fontsize{14}{18}\selectfont}%
    {\makebox[\sechangindent][l]{\thesubsubsection}}{0pt}{}%
    [\global\setlength{\parindent}{\dimexpr2\sechangindent\relax}]%
  \titlespacing{\subsubsection}{\dimexpr2\sechangindent\relax}{8pt}{4pt}%
}

% ─── Table of Contents ───────────────────────────────────────
\usepackage{tocloft}
\renewcommand{\contentsname}{สารบัญ}
% polyglossia เรียก \captionsthai ทุกครั้งที่เปลี่ยนภาษา → ต้อง patch ที่ตรงนี้
\makeatletter
\appto\captionsthai{%
  \def\listtablename{สารบัญตาราง}%
  \def\listfigurename{สารบัญภาพ}%
  \def\figurename{ภาพที่}%
}
\makeatother

% รูปแบบหัวข้อหน้าสารบัญ (18pt หนา จัดกึ่งกลาง% สำหรับสารบัญ (TOC)
\setcounter{tocdepth}{1} % แสดงในสารบัญถึง section (4.1) ซ่อน subsection (4.2.1) ลงไป
\setcounter{secnumdepth}{4} % ยังคงให้มีเลขหมายหัวข้อในเอกสารถึงระดับ 4 (1.1.1.1)

\renewcommand{\cfttoctitlefont}{\hfill\fontsize{18}{22}\selectfont\bfseries}
\renewcommand{\cftaftertoctitle}{\hfill\mbox{}\\[18pt]\hfill หน้า}
\setlength{\cftbeforetoctitleskip}{-\topskip}
\setlength{\cftaftertoctitleskip}{2pt}

\renewcommand{\cftlottitlefont}{\hfill\fontsize{18}{22}\selectfont\bfseries}
\renewcommand{\cftafterlottitle}{\hfill\\[18pt]\hfill หน้า}
\setlength{\cftbeforelottitleskip}{-\topskip}
\setlength{\cftafterlottitleskip}{2pt}

\renewcommand{\cftloftitlefont}{\hfill\fontsize{18}{22}\selectfont\bfseries}
\renewcommand{\cftafterloftitle}{\hfill\\[18pt]\hfill หน้า}
\setlength{\cftbeforeloftitleskip}{-\topskip}
\setlength{\cftafterloftitleskip}{2pt}

% รูปแบบบทในสารบัญ → "บทที่ X  ชื่อบท"
\newlength{\chapprenumwd}
\settowidth{\chapprenumwd}{\bfseries บทที่\ }
\newlength{\secnumwd}
\settowidth{\secnumwd}{5.5\enskip}
\newlength{\subsecnumwd}
\settowidth{\subsecnumwd}{5.5.5\enskip}

\renewcommand{\cftchappresnum}{บทที่\ }
\renewcommand{\cftchapaftersnum}{}
\setlength{\cftchapnumwidth}{\sechangindent}   % ชื่อบทเริ่มตรงเส้นซ้าย 1.27cm

\setlength{\cftsecindent}{\sechangindent}   % ตรงกับย่อหน้า 1.27cm
\setlength{\cftsecnumwidth}{\secnumwd}

\setlength{\cftsubsecindent}{\dimexpr\sechangindent+\secnumwd\relax}
\setlength{\cftsubsecnumwidth}{\subsecnumwd}

\setlength{\cftsubsubsecindent}{\dimexpr\sechangindent+\secnumwd+\subsecnumwd\relax}
\setlength{\cftsubsubsecnumwidth}{1em}

\renewcommand{\cftchapfont}{\mdseries}
\renewcommand{\cftchappagefont}{\mdseries}
\renewcommand{\cftsecfont}{\mdseries}
\renewcommand{\cftsecpagefont}{\mdseries}
\renewcommand{\cftsubsecfont}{\mdseries}
\renewcommand{\cftsubsecpagefont}{\mdseries}
\setlength{\cftbeforechapskip}{0pt}
\renewcommand{\cftchapleader}{\cftdotfill{\cftdotsep}}

% สารบัญตาราง / สารบัญภาพ
\newlength{\tabnumprewd}
\settowidth{\tabnumprewd}{ตารางที่\ }
\newlength{\fignumprewd}
\settowidth{\fignumprewd}{ภาพที่\ }

% ซ่อน addvspace ระหว่างบทใน LOF และ LOT → รายการต่อเนื่องกัน
\makeatletter
\patchcmd{\@chapter}{\addtocontents{lof}{\protect\addvspace{10\p@}}}{}{}{}
\patchcmd{\@chapter}{\addtocontents{lot}{\protect\addvspace{10\p@}}}{}{}{}
\makeatother

\renewcommand{\cfttabpresnum}{ตารางที่\ }
\renewcommand{\cfttabaftersnum}{}
\setlength{\cfttabindent}{0pt}
\setlength{\cfttabnumwidth}{\dimexpr\tabnumprewd+\secnumwd\relax}

\renewcommand{\cftfigpresnum}{ภาพที่\ }
\renewcommand{\cftfigaftersnum}{}
\setlength{\cftfigindent}{0pt}
\setlength{\cftfignumwidth}{\dimexpr\fignumprewd+\secnumwd\relax}

% ─── Figures & Tables ────────────────────────────────────────
\usepackage{graphicx}
\usepackage{float}
\usepackage{booktabs}
\setlength{\fboxsep}{0pt}
\setlength{\fboxrule}{0.5pt}
% รูปภาพในเนื้อหา — ใส่ขอบอัตโนมัติ
\newcommand{\figimg}[2][]{\fbox{\includegraphics[#1]{#2}}}
\usepackage{caption}
\captionsetup[table]{position=top, labelsep=space, font=normalsize, labelfont=bf, singlelinecheck=false, justification=justified, skip=\baselineskip}
\captionsetup[figure]{position=bottom, labelsep=space, font=normalsize, labelfont=bf}
\renewcommand{\tablename}{ตารางที่}
\renewcommand{\figurename}{ภาพที่}

\usepackage{chngcntr}
\counterwithout{figure}{chapter} % ให้เลขรูปภาพรัน 1, 2, 3... ต่อเนื่อง
\counterwithout{table}{chapter}  % ให้เลขตารางรัน 1, 2, 3... ต่อเนื่อง

% ─── Bibliography ────────────────────────────────────────────
\usepackage[
  backend  = biber,
  style    = numeric,
  sorting  = none,
  defernumbers = true
]{biblatex}
\addbibresource{references.bib}

% ─── ปรับรูปแบบบรรณานุกรม → KKU Number Style ─────────────────
% ผลลัพธ์: [N] Author. (Year). Title. Note, จาก URL.
\DeclareFieldFormat{title}{#1}         % ชื่อเรื่อง: ไม่เอียง ไม่มีเครื่องหมายคำพูด
\DeclareFieldFormat{howpublished}{#1}  % URL ใช้ตามที่กำหนด

\DeclareBibliographyDriver{misc}{%
  % ชื่อผู้แต่ง/หน่วยงาน
  \ifnameundef{author}{}{\printnames{author}\addperiod\space}%
  % ปีพิมพ์ในวงเล็บ
  \iffieldundef{year}{}{(\printfield{year})\addperiod\space}%
  % ชื่อเรื่อง
  \iffieldundef{title}{}{\printfield{title}\addperiod\space}%
  % note (เช่น ค้นเมื่อ...)
  \iffieldundef{note}{}{\printfield{note}, \textrm{จาก}\space}%
  % URL
  \iffieldundef{howpublished}{}{\printfield{howpublished}\addperiod}%
  \finentry
}
% ใช้รูปแบบเดียวกันกับ @online
\DeclareBibliographyAlias{online}{misc}


% ─── Miscellaneous ───────────────────────────────────────────
\usepackage{afterpage}
\usepackage{etoolbox}
\usepackage{hyperref}
\hypersetup{
  colorlinks = false,
  hidelinks
}
\usepackage{array}
\usepackage{longtable}
\usepackage{colortbl}
\usepackage{multirow}
\usepackage{pdflscape}
\usepackage{enumitem}
\usepackage{tikz}
\usepackage{xcolor}
\usepackage{listings}
\lstdefinelanguage{json}{
  basicstyle=\ttfamily\small,
  showstringspaces=false,
  breaklines=true,
  frame=single,
  backgroundcolor=\color[gray]{0.97},
  rulecolor=\color[gray]{0.7},
  stringstyle=\color[rgb]{0.13,0.54,0.13},
  morestring=[b]",
}
\lstset{language=json}
\usetikzlibrary{shapes.geometric, shapes.multipart, arrows.meta, positioning, fit, calc, backgrounds, shadows}

% ─── Custom Commands ─────────────────────────────────────────
% ชื่อส่วนประกอบต่าง ๆ (18pt หนา จัดกึ่งกลาง)
\newcommand{\frontsection}[1]{%
  \clearpage
  \thispagestyle{plain}
  \begin{center}
    \fontsize{18}{22}\selectfont\bfseries #1
  \end{center}
  \vspace{12pt}
}
